\section*{Упражнения}
\addcontentsline{toc}{section}{Упражнения}

\begin{exercise}[5-1]
Для $a\in\Q$ обозначим через $G_a$ группу Галуа многочлена $f_a=X^4+X^3+X^2+X+a$. Найдите целые числа
$a_1,a_2,a_3,a_4$ такие, что при $i\ne j$ группы $G_{a_i}$ и $G_{a_j}$ не изоморфны.

\emph{Указание.} При $a=1$ получается круговой многочлен $\Phi_5$ (\rf{5.8}, \rf{5.10}). Для остальных $a$ воспользуйтесь
кубической резольвентой (\rf{4.9}) и дискриминантом:
\[
g_a(X)=X^3-X^2+(1-4a)X+3a-1,\qquad D(f_a)=D(g_a)=256a^3-203a^2+88a-16 .
\]
Выясните, при каких целых $a$ у $g_a$ есть целый корень $r$ (положите $s=4r-3$ и покажите, что $s\mid25$). Для
неприводимых $f_a$ с остальными $a$ таблица после \rf{4.9} оставляет лишь две возможности. Чтобы набрать четыре разные
группы, не ограничивайтесь неприводимыми $f_a$: при некоторых небольших целых $a$ многочлен $f_a$ имеет целый корень
или раскладывается в произведение двух квадратных трёхчленов.
\end{exercise}

\begin{exercise}[5-2]
Докажите, что рациональные решения $a,b\in\Q$ уравнения Пифагора $a^2+b^2=1$ имеют вид
\[
a=\frac{s^2-t^2}{s^2+t^2},\qquad b=\frac{2st}{s^2+t^2},\qquad s,t\in\Q,
\]
и выведите отсюда, что у каждого прямоугольного треугольника с целыми сторонами длины сторон равны
\[
d\,(m^2-n^2,\ 2mn,\ m^2+n^2)
\]
для некоторых целых $d,m,n$.

\emph{Указание.} Условие $a^2+b^2=1$ означает, что $\Nm_{\Q(i)/\Q}(a+bi)=1$ (\rf{5.24}). Примените теорему Гильберта 90
(\rf{5.25}) к $\Q(i)/\Q$ и выразите $\beta/\bar\beta$ через $\beta=s+ti$ (ср.\ \rf{5.25c}). Для второй части возьмите $a=x/z$,
$b=y/z$, где $x,y$~--- катеты, а $z$~--- гипотенуза, сделайте $s,t$ взаимно простыми целыми и разберите отдельно
случаи, когда $s$ и $t$ разной и одинаковой чётности.
\end{exercise}

\begin{exercise}[5-3]
Докажите, что конечное расширение поля $\Q$ содержит лишь конечное число корней из единицы.

\emph{Указание.} Если $E$ содержит первообразный корень $\zeta_n$ степени $n$ из единицы, то $\Q(\zeta_n)\subset E$;
сравните степени с помощью \rf{5.10}. Остаётся показать, что $\varphi(n)\to\infty$: оцените снизу $\varphi(p^k)$ для
степеней простых, делящих $n$.
\end{exercise}

\begin{exercise}[5-4]
Пусть $E$~--- поле разложения неприводимого сепарабельного многочлена $f\in F[X]$. Докажите, что если ни один корень $f$
не порождает $E$ над $F$, то группа $\Gal(E/F)$ содержит ненормальную подгруппу.

\emph{Указание.} Пусть $\alpha$~--- корень $f$ и $H=\Gal(E/F[\alpha])$~--- стабилизатор корня $\alpha$ в $G=\Gal(E/F)$.
По основной теореме теории Галуа (\rf{3.17}) $H$ нормальна в $G$ тогда и только тогда, когда $F[\alpha]=E^H$ нормально над
$F$. Что означала бы нормальность $F[\alpha]/F$ для остальных корней неприводимого многочлена $f$? (Отсюда, в частности:
если группа $\Gal(E/F)$ абелева, то каждый корень $f$ порождает $E$.)
\end{exercise}
